% Build the layer, Chapter 1: the paused holder and the fence.
% Inserted in \S\ref{sec:resource} after the fenced-lease theorem and its
% pitfall, before the ticket lock.
%
% No code or output in this file is typed by hand. \pdexcerpt, \pdrun,
% \pdrunlines and \pdsessionfile print files under excerpts/, which the
% kernel's tools/excerpts.mjs cuts from its source and its recorded runs (see
% figures/pd-listing.tex). To change a listing, change the kernel and cut the
% excerpts again. The prose refers to functions and match arms by name and
% never by line number, so a listing can move within its file.
\subsection{Build the layer: a fence where the write is applied}
\label{sec:swk-build}

\noindent\emph{Express lane: the lease is built twice, as a lock file and then
inside a kernel. One recorded run breaks the first. The same run does not break
the second, and breaks it again when one guard is removed.}

Theorem~\ref{thm:exclusion} describes a gate, and this section builds it. The
code is \pdkernelname{}, a small kernel written for this book. It is not the
reference implementation that the rest of the chapter grades. That daemon has
no such gate; \pdkernelname{} is far smaller and has one.

The section has the same four parts as \weld{prereq:sec:prereq-build}{Naive}:
\emph{Naive}, \emph{Break}, \emph{Fix} and \emph{Verify}. As there, a
script cuts each listing from the kernel's source, and each run is a
recording. The kernel's repository checks both against its source and fails
when they differ.

\subsubsection{Naive: a lock file with a lease}

The usual way to keep two agents out of one file is a lock file beside it. An
agent that wants \texttt{shop/checkout.ts} creates a lock recording its name
and a time. While the lock exists and the time has not passed, the other
agents do not write the file. The time makes the lock a lease, so a holder
that dies does not block everyone forever.

Listing~\ref{alg:swk-build-naive-acquire} is that design, written carefully so
that its failure cannot be blamed on careless code. Creating the lock is one
atomic step: \texttt{tryCreate} writes the record to a private file and
hard-links it into place. The link fails if a lock is already there, so two
agents cannot both create it, and no reader sees half a record. A holder that
acquires again extends its lease. A lock whose time has passed may be broken
by the next agent that tries to acquire it.

\pdexcerpt[caption={The lock file's \texttt{acquire}. Whoever creates the
file holds the lock.},label={alg:swk-build-naive-acquire}]{ch01-naive-acquire}

The second comment admits another race, between two agents that break the
same expired lock. The run below has only one agent breaking the lock, so
that race does not occur.

An agent that uses the lock correctly does one more thing. Just before it
writes, it checks that it still holds the lock and that the lease has not
expired (Listing~\ref{alg:swk-build-check-then-write}). The call to
\texttt{beforeWrite} does nothing in an honest agent. It exists so that a test
can insert a delay at exactly that point.

\pdexcerpt[caption={The careful agent checks its lease, then
writes.},label={alg:swk-build-check-then-write}]{ch01-check-then-write}

The check and the write are two separate actions. The file system applies the
write without consulting the lock, so the result of the agent's check is as
old as the delay between the two. Nothing the agent does removes that gap. A
check can be moved closer to the write, and never into it. \emph{Now you try:}
Exercise~\ref{ex:swk-build-gap} asks for the last point at which the agent's
check can still help.

\subsubsection{Break: the holder pauses}

The scenario is the checkout repair from the paragraph before the theorem.
Heron holds a lease on \texttt{shop/checkout.ts} and intends to raise the tax
rate. Wren is waiting to fix the rounding of the total.

Listing~\ref{alg:swk-build-run-naive} is a run with a pause in the gap. Heron
acquires a lease of 30 seconds, reads the file and computes its edit. It
checks its lease, which is still valid. Then it stalls for 31 seconds, as a
process does during a long garbage collection or a slow network call. The
pause is deterministic. The scenario controls the clock, and it advances the
clock by 31 seconds inside \texttt{beforeWrite}, so every run is the same.

\pdrun[caption={The break. Every step obeys the lock, and Wren's fix is
lost.},label={alg:swk-build-run-naive}]{lost-update-naive}

The times are in the left column. Until \texttt{t=31000} only Heron acts.
By then its lease expired a second ago, so Wren's claim breaks the lock and
is granted. Wren makes its edit under its own lease and releases. Then Heron
resumes. It has already made its check, so it writes. Its edit was computed
from the file as it was before Wren's fix, and it replaces the whole file, so
the tax rate is 0.08 and the rounding is gone.

No step in the run broke a rule. One agent held the lock at a time, and
each agent checked before it wrote. The design fails because a lease can
expire while its holder is paused, and the holder's write carries nothing
that shows how old its permission is. The last lines show a second cost. The
lock file keeps no record, so after the run nothing shows that an update was
lost, or whose it was.

% Landmark 12 asks the question the break has just posed: how can work be
% reassigned without admitting a stale worker? Its page is finished here, so
% that Fix starts on a fresh one. Until this section existed the plate stood
% in the ticket-lock subsection, after that subsection's first paragraph.
\ifdefined\pdswisslandmarkplate
\begin{figure}[p]
  \centering
  \pdswisslandmarkplate{12}{\linewidth}
\end{figure}
\pdmarginpagebreak
\fi

\subsubsection{Fix: number the grants, and check the number at the write}

The fix has two parts, and neither is in the agent. Every grant gets a
number, one more than the previous grant for the same file. And the files
are moved behind one process, the kernel. Every write presents its grant's
number, and the kernel refuses any number except the newest. A holder that
paused past its lease resumes with an old number, and its write is refused.

Listing~\ref{alg:swk-build-claim} issues the numbers. The code calls a number a
\emph{token}; it is the epoch of Theorem~\ref{thm:exclusion}. The function is
passed the newest grant for the path and the time, and it returns a grant or
a reason for refusing. The first arm of the \texttt{match} starts the count
at 1. The second is the one the theorem depends on: a grant that was released
or has expired is followed by a new grant with the next token, whoever claims,
including the old holder. The third renews a live grant for its holder and
keeps the token, so a request that is sent twice gets the same token. The
last refuses any other agent and names the holder.

\pdexcerpt[caption={\texttt{claim}. Each new grant takes the next
token.},label={alg:swk-build-claim}]{ch01-claim}

The second and third arms are in the order that \S\ref{sec:swk-lifecycle}
argued for. The first version of this code had them reversed. A
holder with token 1 that claimed again after its lease had expired then kept
token 1, and its own late writes were admitted. A short model of claims,
written separately from the code, found the mistake. It gave that holder token
2 where the code gave it token 1 again, and the model was right. The Verify
part returns to that check. \emph{Now you try:}
Exercise~\ref{ex:swk-build-arms} asks for a history that distinguishes the two
orders.

Listing~\ref{alg:swk-build-fence} is the check a write must pass. It makes two
checks on the grant that the write presents. \texttt{bound\_to} checks
that the grant was issued to the agent presenting it, since a token is a
small number that anyone could guess. \texttt{newest}, in
Listing~\ref{alg:swk-build-newest}, checks that no later grant exists.
Neither reads the clock. A release does not cause any refusal, and neither
does an expired lease. The next grant does.

\pdexcerpt[caption={\texttt{fence}. The check every write must
pass.},label={alg:swk-build-fence}]{ch01-fence}

\pdexcerpt[caption={\texttt{newest}. An older token is
refused.},label={alg:swk-build-newest}]{ch01-newest}

The first lines of \texttt{newest} are a compile-time switch. With it on,
the function admits every token, and the Verify part uses that build.

These functions only decide. Together they are the kernel's decision core.
They are written in Rust and compiled to WebAssembly, and the module
imports nothing from the program that loads it. So it cannot read a file or a
clock: the time is an argument, like the grants.
Listing~\ref{alg:swk-build-propose} is the code that acts on the decision.
Every change an agent proposes goes through this one method, including a
claim and a release. The method reads the grants for the path and calls the
core. It applies the change if the core admits it, and it logs the decision
in both cases. All of that runs in one database transaction, and the
kernel is the only writer of its database, so no other request runs
between the check and the change.

\pdlistingroom{12}
\pdexcerpt[caption={The kernel's single entry point. A change is checked and
applied in one transaction, and the decision is logged in the same
transaction.},label={alg:swk-build-propose}]{ch01-propose}

\keyidea[Put the check where the write is applied]{A check made by the agent
can be out of date by the time its write arrives. A check made by the process
that applies the write, in the step that applies it, cannot. The kernel adds
a number to each grant and one place that compares it, and that place is the
one that performs the write.}

\subsubsection{Verify: the same run, then the run without the guard}

Listing~\ref{alg:swk-build-run-kernel} is the same run against the kernel. The
agents are the same code and the pause is the same 31 seconds. Only the
workspace they write to has changed.

\pdrun[caption={The same pause against the kernel. The late write is refused,
and the refusal is logged.},label={alg:swk-build-run-kernel}]{lost-update-kernel}

Heron is granted token 1 and Wren token 2. When Heron resumes it presents
token 1. The newest is 2, so the write is refused, and so is the release
after it, for the same reason. The file keeps Wren's rounding. Heron's own
edit was not applied, but Heron has received a refusal, and it can claim the
file again and make the edit on top of Wren's. This time there is also a
record: events 7 and 8 are the two refusals with their reason, written by
the transaction that made them.

A passing run proves less than it appears to. It would also pass if the
write had been refused for some other reason. So the run is repeated against
a kernel with one change: \texttt{newest} is built to admit every token.
Listing~\ref{alg:swk-build-run-no-fence} is the end of that run.

\pdrunlines[caption={The end of the same run with \texttt{newest} switched
off.},label={alg:swk-build-run-no-fence}]{lost-update-no-fence}{1-1,17-32}{1 and 17\textendash 32}

Events 7 and 8 are admitted and the update is lost, as it was under the
lock file. So the guard is what prevented the lost update. The kernel's tests
require both results: the run must pass with the guard and fail without it.
If a guard can be removed without any test failing, the build fails.

The runs have one more check. While they execute, the tests record every
call an agent makes about a claim, with the time it was made and the time it
returned. The calls are a claim, a lease check, a release and a write under
a claim. The tests then search for an order of the calls that is consistent
with those times and in which a sequential model of claims gives every
result the kernel gave. Theorem~\ref{thm:consistency} states that property.
Here it is checked on the histories the tests produce and not proved for
every history. This is the check that found the two arms of \texttt{claim}
in the wrong order.

Table~\ref{tab:swk-build-hypotheses} lists each hypothesis of
Theorem~\ref{thm:exclusion} with the code that satisfies it and the evidence
that the code is what satisfies it.

\begin{table}[H]
\centering
\caption{The hypotheses of Theorem~\ref{thm:exclusion} in \pdkernelname{}:
the code that satisfies each one, and the evidence. [internal]}
\label{tab:swk-build-hypotheses}
\footnotesize
\renewcommand{\arraystretch}{1.2}
\begin{tabularx}{\textwidth}{@{}>{\raggedright\arraybackslash}p{0.31\textwidth} >{\raggedright\arraybackslash}p{0.25\textwidth} >{\raggedright\arraybackslash}X@{}}
\toprule
\textbf{Hypothesis} & \textbf{Code} & \textbf{Evidence} \\
\midrule
Claims on a path are decided one at a time & \texttt{\#propose}: one transaction, one writer & every recorded history fits the sequential model \\
A new grant takes the next token, including a new grant to the old holder & \texttt{claim}, second arm & the model found the arms in the wrong order \\
The count survives a restart & grants are rows in the SQLite file & the kernel is killed after each acknowledged step and restarted; the next claim gets the next token and an old token is still refused. A power loss is not covered (I1b) \\
The process that issues tokens gates the write, in the same step & \texttt{\#propose} & without it: the lock file, Listing~\ref{alg:swk-build-run-naive} \\
The write's token is the newest & \texttt{newest} & without it: Listing~\ref{alg:swk-build-run-no-fence} \\
The token was issued to its presenter & \texttt{bound\_to} & without it: Exercise~\ref{ex:swk-build-unbound} \\
The presenter is the agent it claims to be & a credential issued at \texttt{register} & without the check, an agent that sends another agent's name is accepted, and a test fails \\
\bottomrule
\end{tabularx}
\end{table}

\paragraph{The same kernel, at the command line.}
\pdkernelname{} also runs as a service, and \texttt{harbor} is its
command-line client. The transcript below records two agents and one file.
The service is started with its clock fixed at zero, so that the recording
prints the same times on every run.

\pdsessionfile[two agents and one file, at the command line]{lock-demo}

\texttt{harbor register} asks the service to create an identity. The
service issues a credential and the command stores it. \texttt{HARBOR\_AGENT}
then selects which stored credential a command sends. It does not tell the
service who is calling; the service determines that from the credential
(\S\ref{sec:swk-identity}). Bob's first claim is refused with the holder's
name and the time remaining. After Alice's release his claim is granted with
token 2. The release caused no refusal. The grant that followed it is what
makes token 1 stale.

\begin{pdboundary}
\setlength{\parskip}{4pt plus 1pt}
The runs show that the bad interleaving exists and that the fence refuses it.
They do not show how often an agent pauses past its lease. The fence applies
only to writes that go through the kernel. In \pdkernelname{} the shared
files are rows in the kernel's own database, so its interface is the only
access an agent is given. A process running as the same user can still open
that database, the file \texttt{.harbor/harbord.db}, without going through
the service, which makes this the cooperative rail of
\S\ref{sec:swk-monitor}. Claims are on whole files: the line ranges and
symbols of Table~\ref{tab:swk-claim-granularity} are not built. The costs
named under Theorem~\ref{thm:exclusion} apply here unchanged. A holder whose
lease has expired is still admitted until another agent claims: if Wren had
not claimed, Heron's write at \texttt{t=31000} would have been admitted, a
second after its lease expired. And a power loss can roll the count back.
\end{pdboundary}
